\section{提案モデル}
\label{sec:model}
\subsection{係数生成入力と出力の共通ゲート}
予測対象\(i\)、ラグ\(k\)に対し、入力系列ごとのゲート
\(\bm g_{ki}=(g_{ki1},\ldots,g_{kip})^\top\)と、独立した係数生成器
\(f_{\theta_{ki}}:\R^p\to\R^p\)を設ける。GSE-VARを
\begin{align}
 \bm v_{ki,t}&=\bm g_{ki}\odot\bm x_{t-k},\notag\\
 \bm c_{ki,t}&=\bm g_{ki}\odot f_{\theta_{ki}}(\bm v_{ki,t}),\label{eq:gated-coefficients}\\
 \widehat x_{i,t}&=b_i+\sum_{k=1}^{K}\bm c_{ki,t}^{\top}\bm x_{t-k}
 \label{eq:gse-var}
\end{align}
で定義する。\(\odot\)は要素積、\(b_i\)は入力に依存しない切片、\(\bm c_{ki,t}\)はゲートを含めた有効係数である。ゲートは時点によらない二値変数\(g_{kij}\in\{0,1\}\)とする。

係数生成器は対象・ラグ別の多層パーセプトロン（MLP）とする。同一ラグ内の選択入力間の交互作用を表し、ラグ間は加算で集約する。

\begin{figure}[H]
\centering
\begin{tikzpicture}[>=Stealth,font=\small,
 box/.style={draw,rounded corners=1pt,minimum height=10mm,align=center},
 every node/.style={inner sep=4pt}]
 \node[box,minimum width=16mm] (x) at (0,0) {\(\bm x_{t-k}\)};
 \node[box,minimum width=19mm] (in) at (2.3,0) {入力の選択\\\(\bm v_{ki,t}\)};
 \node[box,minimum width=20mm] (f) at (4.9,0) {係数生成\\\(f_{\theta_{ki}}\)};
 \node[box,minimum width=19mm] (out) at (7.6,0) {出力の選択\\\(\odot\,\bm g_{ki}\)};
 \node[box,minimum width=16mm] (c) at (10,0) {\(\bm c_{ki,t}\)};
 \node[box,fill=gray!8,minimum width=24mm] (g) at (4.9,1.7) {共通の構造ゲート\\\(\bm g_{ki}\)};
 \node[box,minimum width=29mm] (y) at (10,-1.9) {内積・ラグ方向の和\\\(\widehat x_{i,t}\)};
 \draw[->] (x)--(in);
 \draw[->] (in)--(f);
 \draw[->] (f)--(out);
 \draw[->] (out)--(c);
 \draw[->] (g.west)-|(in.north);
 \draw[->] (g.east)-|(out.north);
 \draw[->] (c)--(y);
 \draw[->] (x.south)|-(y.west);
\end{tikzpicture}
\caption{予測対象\(i\)・ラグ\(k\)の計算経路。共通ゲートで係数生成の入力と出力を制御する。下側の過去値の経路も、係数側が\(0\)なら予測に寄与しない。}
\label{fig:architecture}
\end{figure}

\subsection{非因果性の十分条件}
共通ゲートは直接の寄与と係数生成への入力を同時に除く。

\begin{samepage}
\begin{proposition}[入力系列への非依存性]
\label{prop:noncausality}
式\eqref{eq:gated-coefficients}--\eqref{eq:gse-var}のモデルで、固定された\(i,j\)について
\(g_{1ij}=\cdots=g_{Kij}=0\)なら、予測関数\(h_i\)は系列\(j\)の全ラグの値に依存しない。
\end{proposition}
\begin{proof}
\(\bm e_j\)を第\(j\)標準基底ベクトルとし、任意の\(\delta_k\in\R\)によって
\(\widetilde{\bm x}_{t-k}=\bm x_{t-k}+\delta_k\bm e_j\)
と入力を変更する。仮定\(g_{kij}=0\)より、
\begin{equation*}
 \bm g_{ki}\odot\widetilde{\bm x}_{t-k}
 =\bm g_{ki}\odot\bm x_{t-k}
   +\delta_k g_{kij}\bm e_j
 =\bm g_{ki}\odot\bm x_{t-k}.
\end{equation*}
したがって式\eqref{eq:gated-coefficients}で生成される係数は変わらず、
\(\widetilde{\bm c}_{ki,t}=\bm c_{ki,t}\)かつ\(c_{kij,t}=0\)である。
変更後の履歴を\(\widetilde{\bm X}_t\)とすると、式\eqref{eq:gse-var}から
\begin{align*}
 h_i(\widetilde{\bm X}_t)-h_i(\bm X_t)
 &=\sum_{k=1}^{K}\bm c_{ki,t}^{\top}
       (\widetilde{\bm x}_{t-k}-\bm x_{t-k})\\
 &=\sum_{k=1}^{K}\delta_k c_{kij,t}=0.
\end{align*}
\(\delta_k\)は任意なので、予測関数は系列\(j\)の全ラグに依存しない。
\end{proof}
\end{samepage}

この命題は微分可能性によらず、学習後の固定ゲートに成立する。これに基づき、利用可能な入力経路のグラフを
\begin{equation}
 \widehat A_{ij}=\ind\!\left\{\exists k:\ g_{kij}\ne0\right\}
 \label{eq:graph-support}
\end{equation}
と定める。\(\widehat A_{ij}=0\)はモデル内の非依存性を保証し、\(\widehat A_{ij}=1\)は利用可能な経路を表す。後者でも、生成器の出力が\(0\)となる場合や項が相殺される場合には、実際の依存が生じないことがある。

\subsection{有効係数と符号の解釈}
微分可能な点での予測の変化率は
\begin{equation}
 J_{kij,t}
 :=\frac{\partial\widehat x_{i,t}}{\partial x_{j,t-k}}
 =c_{kij,t}
 +\sum_{\ell=1}^{p}x_{\ell,t-k}
       \frac{\partial c_{ki\ell,t}}{\partial x_{j,t-k}}
 \label{eq:jacobian}
\end{equation}
となる。右辺第二項が係数の変化による効果を表すため、係数が一定なら\(J_{kij,t}=c_{kij,t}\)となる。入力依存の場合には両者の対応を分析ごとに確かめることで、重みの方向と限界効果の関係を検討できる。各時点の予測への寄与は、入力値の符号も反映した積\(c_{kij,t}x_{j,t-k}\)で表される。

\subsection{外生入力への拡張}
外生変数を\(\bm u_t\)、利用可能な時点の集合を\(\mathcal K_X\)とする。GSE-VARXでは、内生・外生に別々の係数生成器と二値ゲートを設け、
\begin{equation}
 \widehat x_{i,t}=b_i+\sum_{k=1}^{K}(\bm c^{Y}_{ki,t})^\top\bm x_{t-k}
 +\sum_{k\in\mathcal K_X}(\bm c^{X}_{ki,t})^\top\bm u_{t-k}
\end{equation}
と予測する。各枝の係数は式\eqref{eq:gated-coefficients}と同様に入力・出力の共通ゲートで生成する。外生系列の値は予測対象に含めず、内生・外生の選択確率の和に別々の正則化係数を設定する。既知のカレンダー等には予測時点で利用可能な先行値も指定でき、各入力の系列名と実際のラグを保持する。
